\section{Implementation Details}

\subsection{Backend}
\label{sec:impl-backend}

\verb|job2cool-backend| is a FastAPI application that owns a precise set of \verb|/api|
prefixes (\verb|health|, \verb|chat| for SSE orchestration, \verb|citation|,
\verb|graph_trace|, \verb|score_answer|, \verb|buffers|, \verb|job2cool| for identity,
Projects, and Candidates, \verb|agents|, and \verb|mcp|) and reverse-proxies everything
else to \verb|kb-service|, with \verb|noted| fall-through for PDF files. Seven modules
implement it: \verb|main.py| (app, routes, proxy, Socket.IO mount), \verb|orchestrator.py|
(the agentic pipeline), \verb|services.py| (async LLM and hybrid-RAG clients),
\verb|buffers.py| (live documents), \verb|cache.py| (chunk and turn caches),
\verb|health_monitor.py| (the probe loop), and \verb|socketio_relay.py| (KB-progress and
health pushes).

Live documents are \verb|DocBuffer| objects, in-memory and one per deliverable, with an
\verb|asyncio| pub/sub that broadcasts each change over SSE
(\verb|/api/buffers/events/stream|), so the workspace re-renders the document live, tab by
tab. Token caps are conservative output bounds, not context limits: \verb|SECTION_MAX=8192|
(section composition), \verb|INTRO_MAX=4096| (streamed intro), \verb|DPO_MAX| about 1200
(A2 draft), and \verb|SUMMARY_MAX| about 320 (closing note); \verb|gemma-4|'s context
window is 131,072 tokens, so none of these are system-level constraints. Two FIFO caches in
\verb|cache.py| keep the judge and graph endpoints cheap: a chunk cache (max 2,048, keyed
by \verb|sha1(chunk_id)[:12]|) so citation resolution does not re-query, and a turn cache
(max 256, keyed by \verb|turn_id|) holding the full turn so \verb|/api/score_answer| and
\verb|/api/graph_trace| operate on a finished turn without re-running retrieval. One honest
limitation lives here: \verb|/api/buffers/{id}/save| acknowledges but does not persist to
disk or the KB (a stub), though work still round-trips through the Projects store
(Section~\ref{sec:projects}).

\paragraph{Candidate ingestion.} \verb|scripts/ingest_candidates.py| streams the
\texttt{lang-uk/\allowbreak recruitment-dataset-\allowbreak candidate-profiles-english}
Parquet file (210,250 rows), drops 202 sub-50-character junk profiles, and stores the
remaining 210,048 as \emph{single-chunk} documents. The design originally assumed
section-chunked CVs, but the real CVs are tiny (median about 110 words, or 750 characters),
so one vector per whole CV is both correct and cheaper; the atomic retrieval unit is ``this
candidate,'' and a whole CV fits easily in bge-m3's 8192-token window. The corpus is also
\emph{vector-only}: full graph entity-extraction over 210k CVs benchmarked at about 2.3
days of GPU time and the matcher never reads the graph, so it was deliberately skipped.
Structured fields (\verb|primary_keyword|, \verb|english_level|, \verb|experience_years|,
and \verb|position|) are stored as ChromaDB metadata, the substrate for the
Section~\ref{sec:candidates} pre-filter, and deterministic IDs (\verb|cand-{uuid}|) make
re-runs idempotent.

Honouring the copy-only guardrail, the bulk ingest required two \emph{additive} endpoints
on \verb|noted-rag| rather than edits: \verb|POST /upsert_records| (bulk embed and upsert
with the caller's own flat metadata, token-budgeted sub-batching, \verb|skip_existing|, and
a \textbf{read-back verification} that re-reads the written IDs from ChromaDB before
acknowledging, so a 200 means rows confirmed present) and \verb|POST /list_records|
(paginated browse, powering the Candidates view). The final segment reported
\verb|indexed 43008 (read-back ack 43008)|, every written row DB-confirmed.

The candidate corpus shares a lineage with the A1/A2 training data: both come from the
Djinni platform (job postings for A1/A2, candidate profiles for A3) and share an identical
structured vocabulary (Primary Keyword, English Level, and Experience Years), so a
posting's requirements filter directly into a candidate query with no schema translation.

\subsection{Frontend and UX}

The frontend is a custom single-page application (\verb|frontend/index.html|) baked into
the Docker image at build time, so a frontend or backend change redeploys with a single
\texttt{docker compose up -d -{}-build job2cool-backend}. Layout is a full-viewport flex row: a
left nav rail and a main area with a workspace top bar, a tab bar, and a split document/PDF
pane. Gold (\verb|--primary: #ffe19b|) is the accent; the TalentForge AI mockup is the
\emph{layout} reference, not the colour spec. Each nav view is a small module under
\verb|js2c/|, and Diana is a patched copy of the \verb|cv-chat.js| widget.

The left rail, verified against the running build, is \textbf{Workspace, Projects,
Candidates, Agents, Skills, Tools, Knowledge Base}, and \textbf{Company Profile} (the one
remaining placeholder), with Help \& Support, Settings, and the authenticated-user profile
at the bottom. Diana is patched in five places to live here: the chat POST carries
\verb|JOB2COOL_CONFIG|; a new-turn signal resets tabs; a citation with a source path opens
the PDF split pane (\verb|JOB2COOL_OPEN_PDF|) instead of an in-page highlight; asset paths
are base-relative (\verb|document.baseURI|) for \verb|/job2cool/| sub-path deployment; and
the persona becomes Diana with its own greetings. The widget retains the full stack: the
Thinking, Graph, and Score panels, citations, and STT/TTS/avatar.

The Workspace shows one tab per generated deliverable (created lazily as each section
completes) and an ``All Documents'' overview presenting the deliverables as a cards grid
with a Generation Progress checklist that updates live, the mockup's stepper wired to the
real pipeline. The Job Offer tab shows a Gemma/MA2 segmented toggle when both versions
exist. Per-tab actions are Copy Markdown and PDF export.

The remaining nav views are each a \verb|js2c| module: \verb|kb.js| (Knowledge Base domain
manager with live Socket.IO build progress), \verb|candidates.js| (the Candidates browser),
\verb|agents.js| (CRUD over the \verb|job2cool_orchestrator|, \verb|_composer|, and
\verb|_judge| presets that drive Diana), \verb|mcp.js| (the Skills and Tools admin served
through \verb|mcp-service|), \verb|help.js| and \verb|architecture.js| (the live health
map), \verb|pdfcite.js| (the PDF and bbox renderer), and \verb|sidepanel.js|.

\paragraph{Candidates browser.} \verb|candidates.js| renders a paginated browse over all
210k CVs, an Enter-to-search semantic retrieval whose results carry a match-percentage chip
(the reranker score), and a per-CV detail side-panel. One bug surfaced during the build:
the list page size (30) exceeded \verb|noted-rag|'s \verb|top_k <= 20| cap, so every search
returned an empty 422 until the backend clamped \verb|top_k|; a \verb|setTimeout| search
debounce, ruled out by the project's no-polling rule, was replaced with explicit
Enter-to-search.

\subsection{Projects}
\label{sec:projects}

Conversations are organised into Projects (\verb|chats.js| plus \verb|/api/job2cool/chats/*|),
each with a name, optional description, and \textbf{private or shared visibility (shared by
default)}. Access control is real: private projects live per user; shared projects are
listed for everyone but are owner-only for edit and delete. Reopening a project performs
full-fidelity replay, restoring the conversation, the workspace documents, the role, and
the live Thinking, Graph, and Score panels of every past turn (panel snapshots are
persisted with the project and re-rendered through the same widget code path). This is what
compensates for the buffer-save stub: work persists in the project even though per-buffer
save is deferred.

\subsection{Infrastructure and Live Monitoring}
\label{sec:infra}

job2cool ships as one image (backend plus baked-in frontend) deployed with
\texttt{docker compose up -d -{}-build job2cool-backend} against the shared stack. A
project-wide rule shapes the implementation: \textbf{the browser never polls}.
\verb|setTimeout| and \verb|setInterval| for data-fetching are banned, and live updates are
pushed over SSE for the document buffers and over Socket.IO for KB build progress
(\verb|kb:progress|) and infrastructure health.

The Help \& Support view is a self-documenting, live-monitored architecture map: a
hand-authored, data-driven SVG (no graph library) of sixteen boxes, the browser plus the
fifteen request-path containers, deliberately excluding \verb|noted|'s ingestion and
training infrastructure (Airflow, Postgres, Redis, MLflow, MinIO, and Evidently), which
powers KB operations rather than the request path. Server-side, \verb|health_monitor.py|
runs one \verb|asyncio| loop probing about 13 containers every 10\,s and pushing a
\verb|health:status| snapshot to a Socket.IO room, but only while at least one browser is
watching (watcher-gated). Each probe maps to a colour: a connection error or timeout is
red, an HTTP 5xx is orange, any other response is green. Two cases are special:
\verb|websearch_server|, isolated on \verb|mcp_internal|, is observed through a
side-effect-free passthrough on \verb|mcp-service|; and the backend stamps itself, so if it
dies no push arrives at all, which the client turns into the correct all-red signal.

Detecting the \emph{absence} of an event (staleness) is solved without a timer too. A CSS
animation decays a ``live feed'' badge from green through yellow and orange to red over
about 40\,s, every incoming push restarts the keyframe (an event-driven reset), and if
pushes stop the browser's own animation clock carries the badge to red and the
\verb|animationend| event dims all dots. The monitoring is liveness, not deep health (a
green dot means ``responded,'' not ``every subsystem healthy''), and it is a single-operator
dashboard rather than an alerting pipeline, which is appropriate for a demo and recorded as
such.
